\chapter{Navigation, Autonomy, and Systems Theory}

\section{What navigation means in RUDRAv4 today}

RUDRAv4 currently implements the layers that make navigation possible:

\begin{itemize}
  \item a bounded base command path from PS2 or \topic{/cmd_vel} to the Teensy,
  \item a LiDAR driver publishing \topic{/scan} in the \cmd{laser} frame,
  \item a local obstacle guard that modifies unsafe forward commands,
  \item wheel odometry and IMU telemetry from the Teensy,
  \item an EKF that publishes \topic{/odometry/filtered},
  \item a TF tree that lets RViz and later planners reason about geometry,
  \item DCDC-USB diagnostics for the NUC power rail.
\end{itemize}

It does not yet include a checked-in full Nav2 launch that performs global planning,
local planning, recovery behavior, and goal execution. That is an intentional staging
choice. A navigation stack should be added only after the base can repeatedly prove
that \topic{/scan}, \topic{/odometry/filtered}, \cmd{odom -> base_link}, and the
\topic{/cmd_vel} bridge are stable.

\section{The autonomy stack as layered contracts}

Autonomy is easiest to debug when each layer has a contract:

\begin{longtable}{p{0.22\linewidth}p{0.34\linewidth}p{0.34\linewidth}}
\toprule
Layer & Contract & RUDRAv4 implementation\\
\midrule
Power & Keep the onboard computer alive and observable. & DCDC-USB, \topic{/rudra/nuc_battery}, \topic{/diagnostics}.\\
Actuation & Convert bounded commands into motor driver packets. & Teensy, Sabertooth packet serial, command timeout.\\
Manual authority & Allow a human to command and stop the robot. & Uno PS2 packet stream and \pkg{ps2_uno_to_teensy}.\\
Perception & Publish range data in a known frame. & YDLIDAR G2B on \topic{/scan}.\\
Reactive safety & Reduce or block unsafe forward motion. & \pkg{LidarObstacleGuard}.\\
State estimation & Publish locally continuous robot motion. & wheel odom, IMU yaw rate, \pkg{robot_localization}.\\
Visualization & Make geometry and health visible to the operator. & RViz views, TF, diagnostics.\\
Planning & Turn maps and goals into velocity commands. & future Nav2 or SLAM layer publishing \topic{/cmd_vel}.\\
\bottomrule
\end{longtable}

This structure matters because failures become local. If RViz shows no scan, the
problem is not the Sabertooth. If \topic{/cmd_vel} is correct but the Teensy is silent,
the problem is at the serial bridge or firmware boundary. If the DCDC diagnostics warn
about input sag, the robot may still drive correctly, but the compute rail is no longer
trustworthy enough for long autonomous runs.

\section{Perception: LaserScan as geometry}

A \cmd{LaserScan} is a polar measurement:

\begin{equation}
p_i =
\begin{bmatrix}
r_i\cos\theta_i\\
r_i\sin\theta_i
\end{bmatrix},
\qquad
\theta_i = \theta_{min} + i\Delta\theta.
\end{equation}

The scan is useful only when its frame is known. RUDRAv4 publishes the scan in
\cmd{laser} and uses a static transform from \cmd{base_link} to \cmd{laser}. The
obstacle guard is intentionally simpler than a planner: it examines a front cone,
finds valid finite ranges, and uses the closest obstacle to scale or block forward
motion. This gives immediate local protection without needing a map.

\section{State estimation: odom is not map}

\cmd{odom} is locally continuous. It should not jump. Wheel odometry and IMU yaw rate
are good sources for this frame because they are high-rate and smooth. They drift,
but they drift continuously. A future \cmd{map} frame solves a different problem:
global consistency. SLAM or localization can correct drift by publishing
\cmd{map -> odom}. The separation is:

\begin{equation}
\cmd{map} \rightarrow \cmd{odom} \rightarrow \cmd{base_link}.
\end{equation}

RUDRAv4 should therefore validate \cmd{odom -> base_link} before introducing
\cmd{map -> odom}. If the local frame is unstable, a global localization layer will
hide the problem until the planner behaves unpredictably.

\section{Control: open-loop command, closed-loop roadmap}

The current base is command-bounded but not wheel-speed closed-loop. That means
\topic{/cmd_vel} is converted to left/right Sabertooth command magnitudes, and the
Teensy ramps those commands to reduce abrupt changes. Encoders report what happened,
but they do not yet correct the motor output.

A later PID layer would change the contract. Instead of commanding motor effort
directly, the firmware would command target wheel velocity:

\begin{align}
e(t) &= v_{target}(t)-v_{measured}(t),\\
u(t) &= K_p e(t)+K_i\int e(t)\,dt+K_d\frac{de(t)}{dt}.
\end{align}

The design should add PID only after encoder sign, counts-per-revolution, wheel radius,
track width, and serial timing are all validated. Otherwise the controller will amplify
calibration errors.

\section{Navigation readiness criteria}

Before adding Nav2 or a SLAM launch, RUDRAv4 should pass the following checks:

\begin{enumerate}
  \item \topic{/scan} publishes continuously and appears in RViz with fixed frame \cmd{base_link}.
  \item \topic{/rudra/wheel_odom} changes correctly when the robot moves forward, reverses, and rotates.
  \item \topic{/rudra/imu/raw} reports plausible yaw rate while rotating in place.
  \item \topic{/odometry/filtered} remains stable while stationary and moves smoothly while driving.
  \item TF contains \cmd{odom -> base_link -> laser} and \cmd{base_link -> imu_link}.
  \item The obstacle guard reports \cmd{clear}, \cmd{slowing}, and \cmd{blocked} in predictable physical situations.
  \item The \topic{/cmd_vel} bridge can stop the robot on timeout.
  \item \topic{/diagnostics} reports nominal DCDC output and no sustained input voltage sag.
\end{enumerate}

\section{Operational launch pattern}

The current full robot-side bringup for manual development is the integrated launch:

\begin{lstlisting}[language=bash]
cd /home/rudra/Projects/RUDRAv4
source /opt/ros/lyrical/setup.bash
source /home/rudra/ros2_ws/install/setup.bash
source install/setup.bash
ros2 launch rudra_base_bridge rudra_full.launch.py \
  enable_dcdc_monitor:=true
\end{lstlisting}

For navigation development where voice should not run, keep the same launch and disable voice explicitly:

\begin{lstlisting}[language=bash]
ros2 launch rudra_base_bridge rudra_full.launch.py \
  enable_voice:=false \
  enable_dcdc_monitor:=true
\end{lstlisting}

The current RViz-side view on Aghora is:

\begin{lstlisting}[language=bash]
cd /home/rudra/Projects/RUDRAv4
source /opt/ros/lyrical/setup.bash
export ROS_DOMAIN_ID=42
export ROS_LOCALHOST_ONLY=0
source install/setup.bash
ros2 launch rudra_base_bridge localization_view.launch.py
\end{lstlisting}

For autonomy, stop the PS2 bridge and start the \topic{/cmd_vel} bridge. Then run the
planner or teleop node that publishes \topic{/cmd_vel}. Do not allow two nodes to own
the Teensy serial port at the same time.

\begin{safetybox}{Autonomy gate}
Autonomous navigation should be treated as a privilege earned by the lower layers. If
manual control, scan geometry, TF, odometry, and power diagnostics are not boringly
reliable, a planner will simply move the uncertainty into a faster part of the system.
\end{safetybox}
